\begin{lemma}[A smooth compactly supported function lies in every $W^{1,r}$]
Let $n\ge1$ be an integer and $r\ge1$ real, let $\Omega\subseteq\R^n$ be open,
and let $\varphi\in C^\infty(\R^n)$ have compact support with
$\operatorname{supp}\varphi\subseteq\Omega$.  Then
\[
\varphi\in W^{1,r}(\Omega)\quad\text{with weak gradient}\quad
x\longmapsto\nabla\varphi(x)
\]
(\noderef{MemW1pOn}), where $\nabla\varphi$ is the classical gradient, i.e.\ the
vector field whose $i$-th component is the classical partial derivative
$D_i\varphi$.
\end{lemma}
\begin{proof}
Write $K=\operatorname{supp}\varphi$, a compact subset of $\Omega$.  By
\noderef{MemW1pOn} the assertion has three parts: the test-function identity
\[
\int_\Omega\varphi\,D_i\psi\,dx=-\int_\Omega D_i\varphi\,\psi\,dx
\tag{1}
\]
for every $\psi\in C^\infty(\R^n)$ with compact support contained in $\Omega$
and every $i\in\{1,\dots,n\}$ --- this is \noderef{HasWeakGradientOn} for the
pair $(\varphi,\nabla\varphi)$ --- and the memberships $\varphi\in L^r(\Omega)$,
$|\nabla\varphi|\in L^r(\Omega)$.

\emph{Step 1: $\varphi$ and $\nabla\varphi$ are continuous, bounded, and
supported in $K$.}  Being $C^\infty$, $\varphi$ is continuous and each
$D_i\varphi$ is continuous, so $\nabla\varphi$ is continuous and so is
$x\mapsto|\nabla\varphi(x)|$.  If $x\notin K$ then $\varphi$ vanishes
identically on the open set $\R^n\setminus K$, which contains $x$, hence
$\varphi(x)=0$ and $D\varphi(x)=0$, so also $\nabla\varphi(x)=0$: a function
constant on an open set has vanishing derivative there.  Thus
\[
\varphi=0\ \text{and}\ \nabla\varphi=0\quad\text{on }\R^n\setminus K .
\tag{2}
\]
Since $K$ is compact and $|\varphi|$, $|\nabla\varphi|$ are continuous, they
attain finite maxima on $K$; combined with (2) they are bounded on all of
$\R^n$.  Put
$\Lambda=\sup_{\R^n}|\varphi|<\infty$ and
$\Lambda'=\sup_{\R^n}|\nabla\varphi|<\infty$.

\emph{Step 2: the two $L^r$ memberships.}  Both $\varphi$ and
$|\nabla\varphi|$ are continuous, hence measurable.  By (2) and Step 1,
\[
\int_\Omega|\varphi|^r\,dx=\int_{\Omega\cap K}|\varphi|^r\,dx
\le\Lambda^r\,|K|<\infty ,
\qquad
\int_\Omega|\nabla\varphi|^r\,dx\le(\Lambda')^r\,|K|<\infty ,
\]
using that $|K|<\infty$ because $K$ is compact, and that $t\mapsto t^r$ is
nondecreasing on $[0,\infty)$ so that monotonicity of the integral (Preamble)
applies.  Hence $\varphi\in L^r(\Omega)$ and $|\nabla\varphi|\in L^r(\Omega)$.

\emph{Step 3: the three integrability hypotheses of integration by parts.}  Fix
$\psi\in C^\infty(\R^n)$ with compact support contained in $\Omega$ and fix an
index $i$; write $v=e_i$ for the $i$-th standard basis vector, so that
$D_v=D_i$.  Each of the three functions
\[
x\mapsto D_i\varphi(x)\,\psi(x),
\qquad
x\mapsto \varphi(x)\,D_i\psi(x),
\qquad
x\mapsto \varphi(x)\,\psi(x)
\]
is continuous, being a product of continuous functions, and vanishes off $K$ by
(2) --- for the middle one, $\varphi$ vanishes off $K$; for the other two,
either factor does.  A continuous function vanishing off a compact set is
integrable on $\R^n$: it is bounded, say by $\Lambda''$, so its integral of
absolute values is at most $\Lambda''|K|<\infty$.  So all three are integrable.

\emph{Step 4: integration by parts.}  Both $\varphi$ and $\psi$ are $C^\infty$,
hence differentiable at every point of $\R^n$, so in particular $\varphi$ is
differentiable at every point of $\operatorname{supp}\psi$ and $\psi$ at every
point of $\operatorname{supp}\varphi$.  With Step 3 this is exactly the
hypothesis package of the Preamble item on integration by parts in a direction,
applied with $f=\varphi$, $g=\psi$ and $v=e_i$; it gives
\[
\int_{\R^n}\varphi\,D_i\psi\,dx=-\int_{\R^n}D_i\varphi\,\psi\,dx .
\tag{3}
\]

\emph{Step 5: from $\R^n$ to $\Omega$.}  Both integrands in (3) vanish at every
point outside $K$, by (2); and $K\subseteq\Omega$.  Hence for each of them the
integral over $\Omega$ and the integral over $\R^n$ agree: the difference is the
integral over $\R^n\setminus\Omega$ of a function that is identically $0$ there,
since $\R^n\setminus\Omega\subseteq\R^n\setminus K$.  Therefore (3) reads
\[
\int_\Omega\varphi\,D_i\psi\,dx=-\int_\Omega D_i\varphi\,\psi\,dx ,
\]
which is (1); and since $D_i\varphi$ is the $i$-th component of
$\nabla\varphi$, this is the defining identity of
\noderef{HasWeakGradientOn} for the pair $(\varphi,\nabla\varphi)$ on $\Omega$
at the index $i$.  As $\psi$ and $i$ were arbitrary, that identity holds for
the whole test class.

Together with Step 2 this is the assertion $\varphi\in W^{1,r}(\Omega)$ with
weak gradient $\nabla\varphi$ (\noderef{MemW1pOn}).  Note that the exponent $r$
entered only in Step 2, through the bound $\Lambda^r|K|<\infty$, so the
conclusion holds simultaneously for every $r\ge1$.
\end{proof}
